\documentclass[11pt]{article}
\usepackage[T1]{fontenc}
\usepackage{lmodern}
\usepackage{amsmath,amssymb,amsthm,mathtools}
\usepackage{graphicx}
\usepackage[margin=1in]{geometry}
\usepackage[colorlinks=true,linkcolor=blue,citecolor=blue,urlcolor=blue]{hyperref}
\usepackage{microtype}

\newtheorem{theorem}{Theorem}
\newtheorem{lemma}[theorem]{Lemma}
\newtheorem{remark}[theorem]{Remark}
\newcommand{\N}{\mathbb N}
\newcommand{\ellone}{\ell_1}
\newcommand{\snorm}[1]{\lVert #1\rVert_s}
\newcommand{\onorm}[1]{\lVert #1\rVert_1}

\title{Fixed-branching diamonds factor through the summing basis\\
\large A positive answer to Remark 1.7 of arXiv:1609.06618}
\author{Automated mathematics discovery run \texttt{fa\_banach\_001}\\
Lane 12, candidate for human review}
\date{21 July 2026}

\begin{document}
\maketitle

\begin{center}
\fbox{\parbox{0.92\textwidth}{\textbf{Status: candidate full solution, likely
valid.}  For every fixed finite branching number $k$, we construct the
factorization maps asked for in Remark 1.7.  If $C_2$ is a binary-diamond
constant, the construction gives $C(k)=2C_2\lceil\log_2 k\rceil$.  The new
ingredient is an elementary family of binary projections, one of which
preserves each prescribed pairwise distance exactly.}}
\end{center}

\section{Source question}

Ostrovskii and Randrianantoanina study multibranching diamonds
$D_{n,k}$ and factorization between the unit-vector and summing bases
\cite{OR}.  Their condition (1.1) asks for maps $f_n:D_{n,k}\to\ellone$ such
that, uniformly in $n$,
\begin{equation}\label{eq:factor}
 \onorm{f_n(u)-f_n(v)}\le d_{D_{n,k}}(u,v)
 < C(k)\snorm{f_n(u)-f_n(v)}.
\end{equation}
Here
\[
 \snorm{(a_i)}=\sup_r\left|\sum_{i=1}^r a_i\right|.
\]
On page 4, Remark 1.7 asks whether $C(k)$ exists for every fixed $k$ and
emphasizes that even $C(3)$ was unknown:

\begin{center}
\includegraphics[width=0.98\textwidth]{figures/open_question_crop.png}
\end{center}

We answer the entire fixed-branching question affirmatively.

\section{Statement of the result}

We use the recursive description of $D_{n,k}$.  The graph $D_{0,k}$ is one
edge with bottom and top endpoints.  To obtain $D_{n,k}$, replace every edge
by $k$ paths of length two, with the edge lengths halved.  Equivalently,
$D_{n,k}$ is the parallel composition of $k$ arms, each arm being the series
composition of two copies of $D_{n-1,k}$ scaled by $1/2$.

Johnson and Schechtman proved the binary case \cite{JS}.  Thus there is a
constant $C_2>1$ such that, for every $n$, some
$g_n:D_{n,2}\to\ellone$ satisfies
\begin{equation}\label{eq:binary}
 \onorm{g_n(a)-g_n(b)}\le d_{D_{n,2}}(a,b)
 < C_2\snorm{g_n(a)-g_n(b)}.
\end{equation}

\begin{theorem}[All fixed branchings]\label{thm:main}
For every integer $k\ge2$ and every $n\in\N$, there is a map
$F_{n,k}:D_{n,k}\to\ellone$ satisfying
\[
 \onorm{F_{n,k}(x)-F_{n,k}(y)}\le d_{D_{n,k}}(x,y)
 < 2C_2\lceil\log_2 k\rceil
   \snorm{F_{n,k}(x)-F_{n,k}(y)}
\]
for all distinct vertices $x,y$.  In particular, the numbers $C(k)$ asked
for in Remark 1.7 exist for every finite $k$, and one may take
\[
                 C(k)=2C_2\lceil\log_2 k\rceil.
\]
\end{theorem}

The logarithmic upper bound is compatible with Theorem 1.6 of \cite{OR},
which implies that valid constants cannot remain bounded as $k\to\infty$.
We make no claim that the logarithmic order or its coefficient is sharp.

\section{Idea of the proof}

Assign distinct binary words to the $k$ arm labels.  Reading one coordinate
of those words collapses every $k$-fold parallel replacement onto a binary
replacement.  The resulting map to a binary diamond preserves the height of
every vertex in every recursive cell.  A pair of vertices has at most one
relevant parallel split: the first recursive cell in which the two points
occupy different arms.  Any bit distinguishing those two arm labels preserves
the pair's distance exactly.

We therefore have $m=\lceil\log_2 k\rceil$ nonexpansive binary images and at
least one exact image for every pair.  Compose all images with $g_n$, put them
in consecutive coordinate blocks, and average by $1/m$.  The $\ellone$ upper
bound averages.  An arbitrary offset from earlier blocks can reduce the
summing norm visible inside the exact block by at most a factor two.  This
gives $2mC_2$.

\section{Binary projections}

Fix a function $q:\{1,\ldots,k\}\to\{0,1\}$.  It induces recursively a map
\[
                    \pi_q:D_{n,k}\longrightarrow D_{n,2}.
\]
The bottom and top go to the binary bottom and top.  In every replaced edge,
the midpoint in source arm $i$ goes to the midpoint in binary arm $q(i)$;
inside each of the two child edges the same rule is repeated.  Multiple
source arms are allowed to collapse onto one binary arm.

\begin{lemma}[One binary projection preserves each distance]
\label{lem:projection}
Every $\pi_q$ is $1$-Lipschitz and preserves the distance from a vertex to
the bottom and top of every recursive subdiamond containing that vertex.
Moreover, for every $x,y\in D_{n,k}$, one of the following holds:
\begin{enumerate}
 \item every map $\pi_q$ preserves $d(x,y)$; or
 \item there are two distinct arm labels $a,b\in\{1,\ldots,k\}$ such that
       every $q$ with $q(a)\ne q(b)$ satisfies
       \[
          d_{D_{n,2}}(\pi_q(x),\pi_q(y))=d_{D_{n,k}}(x,y).
       \]
\end{enumerate}
\end{lemma}

\begin{proof}
We argue by induction on $n$, simultaneously for the endpoint assertion and
the dichotomy.  The result is immediate for $n=0$.  A final edge of
$D_{n,k}$ maps to a final edge of $D_{n,2}$ of the same length, so $\pi_q$ is
$1$-Lipschitz.

Let the bottom and top of the root diamond be $B,T$, and let $M_i$ be the
midpoint of arm $i$.  Each lower piece $[B,M_i]$ and upper piece $[M_i,T]$
is a scaled copy of $D_{n-1,k}$.  The induction hypothesis in these pieces
shows that distance to their endpoints is preserved.  Consequently the
distance from any vertex to $B$ and $T$ is preserved as well.  For example,
if $x$ is in the upper piece of arm $i$, the route from $B$ through $M_i$ has
length
\[
             d(B,M_i)+d(M_i,x),
\]
and no route entering through $T$ is shorter; the same statement holds after
projection.  The other cases are symmetric.  Applying the same reasoning
inside any recursive cell proves the full endpoint assertion.

Now consider $x,y$.  Cases in which one point is $B$ or $T$ are already
covered by endpoint preservation.  If $x$ and $y$ lie in distinct root arms
$i\ne j$, put
\[
 h_x=d(B,x),\qquad h_y=d(B,y),\qquad L=d(B,T).
\]
Different arms meet only at $B$ and $T$, hence
\begin{equation}\label{eq:different-arms}
 d(x,y)=\min\{h_x+h_y,\,2L-h_x-h_y\}.
\end{equation}
If $q(i)\ne q(j)$, their images lie in distinct binary arms.  Endpoint
preservation gives the same $h_x,h_y,L$, so the binary analogue of
\eqref{eq:different-arms} is identical.  Alternative (2) holds with
$a=i,b=j$.

It remains to suppose that both points lie in one root arm $i$.  If they lie
in the same lower piece or the same upper piece, apply the inductive
dichotomy in that copy of $D_{n-1,k}$.  If they lie in different pieces, say
$x\in[B,M_i]$ and $y\in[M_i,T]$, every path either passes through $M_i$ or
exits through both $B$ and $T$.  The former route is no longer than the latter:
writing $\alpha=d(B,x)$ and $\beta=d(y,T)$, their lengths are respectively
\[
             L-\alpha-\beta
       \quad\hbox{and}\quad
             L+\alpha+\beta.
\]
Thus
$d(x,y)=d(x,M_i)+d(M_i,y)$.  Both terms are endpoint distances inside child
diamonds and are preserved by every $\pi_q$.  This is alternative (1) and
completes the induction.
\end{proof}

Choose pairwise distinct words
\[
       \sigma(1),\ldots,\sigma(k)\in\{0,1\}^m,
       \qquad m=\lceil\log_2 k\rceil,
\]
and let $q_j(i)=\sigma(i)_j$.  Lemma~\ref{lem:projection} immediately gives
the key metric identity
\begin{equation}\label{eq:max-projection}
 d_{D_{n,k}}(x,y)
 =\max_{1\le j\le m}
   d_{D_{n,2}}(\pi_{q_j}(x),\pi_{q_j}(y)).
\end{equation}
Indeed, in alternative (2), the distinct words of $a$ and $b$ differ in some
coordinate; alternative (1) is even stronger.  The reverse inequality in
\eqref{eq:max-projection} follows from $1$-Lipschitzness.

\section{The summing norm of concatenated blocks}

For finite sequences $z^{(1)},\ldots,z^{(m)}$, write
$z=z^{(1)}\mathbin\Vert\cdots\mathbin\Vert z^{(m)}$ for their concatenation.

\begin{lemma}[Block visibility]\label{lem:block}
For every block $j$,
\[
                         \snorm{z}\ge\frac12\snorm{z^{(j)}}.
\]
\end{lemma}

\begin{proof}
Let $A$ be the sum of all coordinates preceding block $j$, and let $P_r$ be
the partial sums within that block, with $P_0=0$.  Choose $r$ with
$|P_r|=\snorm{z^{(j)}}=:s$.  The global partial sums include both $A$ and
$A+P_r$.  Since
\[
                  s=|P_r|\le |A|+|A+P_r|,
\]
one of those two global partial sums has absolute value at least $s/2$.
\end{proof}

\section{Proof of Theorem~\ref{thm:main}}

For $x\in D_{n,k}$ define
\begin{equation}\label{eq:construction}
 F_{n,k}(x)=\frac1m\bigl(
   g_n(\pi_{q_1}(x))\mathbin\Vert\cdots\mathbin\Vert
   g_n(\pi_{q_m}(x))\bigr).
\end{equation}
Because $g_n$ and every $\pi_{q_j}$ are nonexpansive in the relevant
$\ellone$ and graph metrics,
\begin{align*}
 \onorm{F_{n,k}(x)-F_{n,k}(y)}
 &=\frac1m\sum_{j=1}^m
   \onorm{g_n(\pi_{q_j}(x))-g_n(\pi_{q_j}(y))}\\
 &\le\frac1m\sum_{j=1}^m
   d(\pi_{q_j}(x),\pi_{q_j}(y))
 \le d(x,y).
\end{align*}

By \eqref{eq:max-projection}, choose $j$ for which
$d(\pi_{q_j}(x),\pi_{q_j}(y))=d(x,y)$.  Lemma~\ref{lem:block} and the binary
estimate \eqref{eq:binary} give
\begin{align*}
 \snorm{F_{n,k}(x)-F_{n,k}(y)}
 &\ge \frac{1}{2m}
 \snorm{g_n(\pi_{q_j}(x))-g_n(\pi_{q_j}(y))}\\
 &>\frac{d(x,y)}{2mC_2}.
\end{align*}
This is the claimed strict inverse inequality and completes the proof.

\section{Verification, scope, and novelty}

The packet includes an independent address-based exhaustive checker.  It
constructs all bit projections, computes all-pairs shortest-path distances,
and verifies \eqref{eq:max-projection} at finite depths before checking the
concatenated Hamming realization.  This is only a regression test; the
inductive proof above covers all $n$.

The result answers exactly the full fixed-branching question in Remark 1.7,
including the highlighted case $k=3$.  It does not solve the neighboring
Planar Conjecture, forbidden-minor conjecture, or the problem of embedding all
series-parallel metrics into every non-superreflexive Banach space.

Before promotion, the run's solution, attempt, proof-gap and registry indexes
were searched by arXiv id, exact question, and the phrases
``multibranching diamond'', ``summing norm'', ``$D_{n,3}$'', and ``$C(3)$''.
A bounded primary arXiv search used the same phrases together with ``binary
projection'' and ``fixed branching''.  It found no later resolution or this
binary-code construction.  This is not an exhaustive novelty certification;
the separate \texttt{novelty.md} records the search bounds.

Human review is strongly recommended.  The principal audit target is
Lemma~\ref{lem:projection}, especially the same-arm/different-half case.  The
new proof otherwise uses only the published binary case and elementary norm
estimates.

\begin{thebibliography}{9}
\bibitem{OR}
M. I. Ostrovskii and B. Randrianantoanina,
\emph{A new approach to low-distortion embeddings of finite metric spaces
into non-superreflexive Banach spaces}, arXiv:1609.06618 (2016), Remark 1.7.

\bibitem{JS}
W. B. Johnson and G. Schechtman,
\emph{Diamond graphs and super-reflexivity},
J. Topol. Anal. \textbf{1} (2009), 177--189,
doi:10.1142/S1793525309000114.
\end{thebibliography}

\end{document}
